% TeX
\chapter{本文中英文对照表}

表~\ref{tab:appendix-bilingual} 给出正文常用术语的中英文对照；第二列在必要时附缩写或原文记号。

\begingroup
\small
\setlength{\LTpre}{6pt}
\setlength{\LTpost}{6pt}
\renewcommand{\arraystretch}{1.12}
\begin{longtable}{@{}p{5.6cm}p{7.4cm}@{}}
	\caption{本文主要术语中英文对照表}\label{tab:appendix-bilingual} \\
	\toprule
	\textbf{中文} & \textbf{英文} \\
	\midrule
	\endfirsthead
	\multicolumn{2}{@{}l@{}}{\small 续表~\thetable} \\
	\toprule
	\textbf{中文} & \textbf{英文} \\
	\midrule
	\endhead
	\midrule
	\multicolumn{2}{r@{}}{\small 续下页} \\
	\endfoot
	\bottomrule
	\endlastfoot
	% --- 第 1 章：基础概念与模型 ---
	自注意力机制 & Self-Attention \\
	深度神经网络 & Deep Neural Network (DNN) \\
	卷积神经网络 & Convolutional Neural Network (CNN) \\
	循环神经网络 & Recurrent Neural Network (RNN) \\
	计算机视觉 & Computer Vision (CV) \\
	自然语言处理 & Natural Language Processing (NLP) \\
	语音识别 & Speech Recognition (SR) \\
	大语言模型 & Large Language Model (LLM) \\
	多模态大语言模型 & Multimodal Large Language Model (MLLM) \\
	通用人工智能 & Artificial General Intelligence (AGI) \\
	迁移学习 & Transfer Learning \\
	少样本学习 & Few-Shot Learning \\
	零样本学习 & Zero-Shot Learning \\
	扩展定律 & Scaling Laws \\
	图文对比学习CLIP & Contrastive Language-Image Pre-Training\\
	
	最优先进水平 & state-of-the-art (SOTA) \\
	% --- 第 1--2 章：MLLM 结构 ---
	模态编码器 & Modality Encoder \\
	输入投影层 & Input Projector \\
	LLM 基座 & LLM Backbone \\
	输出投影层 & Output Projector \\
	模态生成器 & Modality Generator \\
	扩散模型 & diffusion model \\
	模块异构 & module heterogeneity \\
	% --- 第 2 章：训练阶段与数据 ---
	模态对齐 & modality alignment \\
	指令微调 & instruction tuning \\
	% --- 并行、通信与显存 ---
	数据并行 & data parallelism (DP) \\
	分布式数据并行 & Distributed Data Parallel (DDP) \\
	张量并行 & tensor parallelism (TP) \\
	流水线并行 & pipeline parallelism (PP) \\
	序列并行 & sequence parallelism (SP) \\
	上下文并行 & context parallelism (CP) \\
	专家并行 & expert parallelism (EP) \\
	混合专家 & Mixture of Experts (MoE) \\
	混合并行 & hybrid parallelism\\
	微批次 & micro-batch \\
	小批量 & mini-batch \\
	一前向一反向调度 & One-Forward-One-Backward (1F1B) \\
	全归约 & AllReduce \\
	全收集 & All-Gather \\
	归约分散 & Reduce-Scatter \\
	全对全通信 & All-to-All \\
	点对点通信 & Peer-to-Peer (P2P) \\
	参数服务器 & parameter server \\
	流水级 & pipeline stage \\
	流水线气泡 & pipeline bubble \\
	混合精度训练 & mixed-precision training \\
	全分片数据并行 & Fully Sharded Data Parallel (FSDP) \\
	零冗余优化器 & Zero Redundancy Optimizer (ZeRO) \\
	环形注意力 & Ring Attention \\
	激活/激活值 & activation \\
	显存溢出 & Out Of Memory (OOM) \\
	负载均衡 & load balancing \\
	剖析/性能剖析 & profiling \\
	设备拓扑 & device topology \\
	整数规划 & integer programming \\
	复合粒度（模型抽象） & multi-granularity model abstraction \\
	成本模型 & cost model \\
	服务等级目标 & Service Level Objective (SLO) \\
	模型算力利用率 & Model FLOPs Utilization (MFU) \\
	% --- 框架与代表性系统（第 1--2 章） ---
	虚拟工作节点 & Virtual Worker (VW) \\
	剖析器 & Profiler\\
	GPU 分配器 & GPU Allocator \\
	阶段放置器 & Stage Placer \\
	全部重计算 & full recomputation \\
	选择重计算 & selective recomputation \\
	细粒度重计算 & fine-grained recomputation \\
	Top-$k$专家路由 & Top-$k$ (expert routing) \\
\end{longtable}
\endgroup
